% GENERATED FILE. Edit canonical lessons, then run npm run generate:books.
% canonical-source-hash: fnv1a64:71bade27c01859dd
% canonical-lessons: PT-C44-carne, PT-C44-peixe, PT-C44-manteiga, PT-C44-sal, PT-C44-frango

\chapter{Meat, Fish, Butter, Salt, Chicken}
\label{ch:pt-meat-fish-butter-salt-chicken}
\hlchaptermodality{44}

\begin{chapteropening}
\textbf{By the end of this chapter:} \emph{I can name meat, fish, butter, salt and chicken.}

Complete the last lesson of chapter 44: I can name meat, fish, butter, salt and chicken.
\end{chapteropening}

\section[a carne]{a carne — meat}

\label{lesson:PT-C44-carne}



\begin{quote}
Before the new one: say the Portuguese for to give, then the Portuguese for to live.
\end{quote}

\subsection*{You'll want to know: a carne}
\textbf{a carne} — \textquotedblleft{}meat\textquotedblright{}.

\begin{cousinweb}
\textbf{carne} $\leftarrow$ Latin \emph{caro, carnis}, \textquotedblleft{}flesh\textquotedblright{} (EtymDB 2.1) — inherited.
\end{cousinweb}

\begin{grammarlens}[title={What you've built}]
One more word for this chapter.
\end{grammarlens}

\subsection*{Your turn}
\begin{itemize}
\raggedright
  \item \emph{Say it:} \emph{a carne}
  \item \emph{Say it:} \emph{a carne}, once more
  \item \emph{Say it:} say \emph{a carne}
  \item \emph{Recall:} say the Portuguese for to give, then the Portuguese for to live, then say \emph{a carne} again
\end{itemize}

\begin{tcolorbox}[breakable,colback=teal!4,colframe=teal!35!black,title={Before you move on}]
What does a carne mean? (\textquotedblleft{}Meat\textquotedblright{}.) Say it once more.
\end{tcolorbox}
\section[o peixe]{o peixe — fish}

\label{lesson:PT-C44-peixe}



\begin{quote}
Before the new one: say the Portuguese for to live, then the Portuguese for meat.
\end{quote}

\subsection*{You'll want to know: o peixe}
\textbf{o peixe} — \textquotedblleft{}fish\textquotedblright{}.

\begin{cousinweb}
\textbf{peixe} $\leftarrow$ Latin \emph{piscis}, \textquotedblleft{}fish\textquotedblright{} — inherited. EtymDB 2.1 traces it only to Old Portuguese \emph{peixe}; it derives Italian \emph{pesce} and Spanish \emph{pez} from \emph{piscis}.
\end{cousinweb}

\begin{grammarlens}[title={What you've built}]
One more word for this chapter.
\end{grammarlens}

\subsection*{Your turn}
\begin{itemize}
\raggedright
  \item \emph{Say it:} \emph{o peixe}
  \item \emph{Say it:} \emph{o peixe}, once more
  \item \emph{Say it:} say \emph{o peixe}
  \item \emph{Recall:} say the Portuguese for to live, then the Portuguese for meat, then say \emph{o peixe} again
\end{itemize}

\begin{tcolorbox}[breakable,colback=teal!4,colframe=teal!35!black,title={Before you move on}]
What does o peixe mean? (\textquotedblleft{}Fish\textquotedblright{}.) Say it once more.
\end{tcolorbox}
\section[a manteiga]{a manteiga — butter}

\label{lesson:PT-C44-manteiga}



\begin{quote}
Before the new one: say the Portuguese for meat, then the Portuguese for fish.
\end{quote}

\subsection*{You'll want to know: a manteiga}
\textbf{a manteiga} — \textquotedblleft{}butter\textquotedblright{}.

\begin{cousinweb}
\textbf{manteiga} — its origin is uncertain, perhaps a word from a language spoken in Iberia before Latin. EtymDB 2.1 lists only cognates, such as Spanish \emph{manteca}, \textquotedblleft{}lard\textquotedblright{}.
\end{cousinweb}

\begin{grammarlens}[title={What you've built}]
One more word for this chapter.
\end{grammarlens}

\subsection*{Your turn}
\begin{itemize}
\raggedright
  \item \emph{Say it:} \emph{a manteiga}
  \item \emph{Say it:} \emph{a manteiga}, once more
  \item \emph{Say it:} say \emph{a manteiga}
  \item \emph{Recall:} say the Portuguese for meat, then the Portuguese for fish, then say \emph{a manteiga} again
\end{itemize}

\begin{tcolorbox}[breakable,colback=teal!4,colframe=teal!35!black,title={Before you move on}]
What does a manteiga mean? (\textquotedblleft{}Butter\textquotedblright{}.) Say it once more.
\end{tcolorbox}
\section[o sal]{o sal — salt}

\label{lesson:PT-C44-sal}



\begin{quote}
Before the new one: say the Portuguese for fish, then the Portuguese for butter.
\end{quote}

\subsection*{You'll want to know: o sal}
\textbf{o sal} — \textquotedblleft{}salt\textquotedblright{}.

\begin{cousinweb}
\textbf{sal} $\leftarrow$ Latin \emph{sal}, \textquotedblleft{}salt\textquotedblright{} (EtymDB 2.1) — inherited.
\end{cousinweb}

\begin{grammarlens}[title={What you've built}]
One more word for this chapter.
\end{grammarlens}

\subsection*{Your turn}
\begin{itemize}
\raggedright
  \item \emph{Say it:} \emph{o sal}
  \item \emph{Say it:} \emph{o sal}, once more
  \item \emph{Say it:} say \emph{o sal}
  \item \emph{Recall:} say the Portuguese for fish, then the Portuguese for butter, then say \emph{o sal} again
\end{itemize}

\begin{tcolorbox}[breakable,colback=teal!4,colframe=teal!35!black,title={Before you move on}]
What does o sal mean? (\textquotedblleft{}Salt\textquotedblright{}.) Say it once more.
\end{tcolorbox}
\section[o frango]{o frango — chicken (to eat)}

\label{lesson:PT-C44-frango}



\begin{quote}
Before the new one: say the Portuguese for butter, then the Portuguese for salt.
\end{quote}

\subsection*{You'll want to know: o frango}
\textbf{o frango} — \textquotedblleft{}chicken (to eat)\textquotedblright{}.

\begin{cousinweb}
\textbf{frango} — its origin is uncertain, and EtymDB 2.1 records no ancestor for it.
\end{cousinweb}

\begin{grammarlens}[title={What you've built}]
One more word for this chapter.
\end{grammarlens}

\subsection*{Your turn}
\begin{itemize}
\raggedright
  \item \emph{Say it:} \emph{o frango}
  \item \emph{Say it:} \emph{o frango}, once more
  \item \emph{Say it:} say \emph{o frango}
  \item \emph{Recall:} say the Portuguese for butter, then the Portuguese for salt, then say \emph{o frango} again
\end{itemize}

\begin{tcolorbox}[breakable,colback=teal!4,colframe=teal!35!black,title={Before you move on}]
What does o frango mean? (\textquotedblleft{}Chicken (to eat)\textquotedblright{}.) Say it once more.
\end{tcolorbox}
